\begin{abstract}
Tree speculative decoding for recurrent-hybrid language models must preserve both candidate ancestry and the state of the selected continuation. We present a systems study of GDN Tree-Scan, a Gated DeltaNet verifier integrated with vLLM and FlashAttention-2. The implementation coordinates branch-local recurrent computation, convolution history, attention key--value remapping, and device-side selection. Archived regression evidence identifies two distinct failures: incomplete accepted-path state publication and a physical-layout change that perturbs the shared attention spine despite masking. The publication repair changes a token fixture from 15 failures in 15 checks to zero; the layout repair produces byte agreement on the shared spine in its tested route. Performance does not follow acceptance alone. In a historical batch-four-configured campaign on Qwen3.6-27B-FP8 and GB10, a tree commits 5.286 tokens per request verification event against a five-step native chain's 4.422, but its complete decode-step rate is 32.85 versus 42.74 tokens/s. A separate paired attention optimization reduces its tested step by 27.03 milliseconds. A kernel-level pilot on identical operands finds that the sequential scan, a forward-substitution (WY-family) solver, and a finite-Neumann (Bole-family) solver all meet a float64 oracle at the fp32 floor when dense products run in fp32, that tf32 products cost three orders of magnitude, and that the production scan reproduces the native verify kernel bit-for-bit on the tested operands in an fp32-store diagnostic; it does not test model-level continuation or fresh prefixes. These findings are scoped to archived configurations, with effective sampling provenance and a matched final-stack comparison still pending. The study provides concrete state and numerical witnesses and reconciled cost accounting, rather than a general superiority claim or a deployed distribution-equivalence proof.
\end{abstract}
